\section*{Estructuras VHDL}

Un proyecto en VHDL se estructura de la siguiente forma.

Un proyecto en VHDL se estructura en 3 partes: 
\begin{enumerate}
    \item Importar las librerías.
    
        Todos los proyectos en VHDL incoran diferentes librerías para poder
        utilizar los tipos.

    \item Definición de puertos y genéricos.
    
        En este apartado se definen dos cosas: los génericos y los puertos.
        Los genéricos son constantes utilizadas para escalar los módulos en VHDL.
        \begin{quotation}
            \noindent Los genéricos en muchos casos son utilizados como constantes generales como en 
            otros lenguajes.
            Esto es un problema que a la larga genera problemas de escalabilidad.
        \end{quotation}

    \item Definición de la arquitectura.
    
        Dentro de la arquitectura se define la operatibidad del módulo FW
    

\end{enumerate}



\begin{verbatim}
    library ieee;
    use ieee.std_logic_1164.all;

    entity modulo is
        generic(

        );
        port(

        );
    end entity;

    architecture arch_modulo of modulo is

    begin

    end architecture;
\end{verbatim}

\subsection{Librerías}
Las librarías en VHDL incorporan los elementos necesarios para poder utilizar el FW.
    \begin{verbatim}
        librarie ieee;
        use ieee.std_logic_1164.all;
    \end{verbatim}
Para incluir librerías propias se incorporan utilizando la librería \textit{work}.
    \begin{verbatim}
        use work.package.all;
    \end{verbatim}

Los fabricantes de FPGAs o de ASICs suelen crear sus propias
librerías para acceder a recursos propios, 
como por ejemplo la librería \textit{UNISIM} de Xilinx(AMD).


\subsection{Entity}
En la entidad se compone de dos partes.

\subsubsection{Generic}


    \begin{verbatim}
        generic(
            G_WIDTH  : integer := 8;
            G_DATA  : integer range 0 to 10 := 8;
            G_TEXT : string := "YES"
        );
    \end{verbatim}

\subsubsection{Port}
En este apartado se definen los puertos del módulo que tienen señales de entrada, 
salida o bidereccionales.

Entre estos puertos suelen incorporar puertos de reloj, de reset o de habilitación.

    \begin{quote}
        Los puertos de reset de normal son activos a nivel bajo ('0').
    \end{quote}

    \begin{verbatim}
        port(
            PUERTO_A_I  : in std_logic;
            PUERTO_B_O  : out std_logic_vector(G_WIDTH-1 downto 0);
            PUERTO_C_IO : inout std_logic
        );
    \end{verbatim}



\subsection{Architecture}

La arquitectura define el comportamiento interno del FW.